\section{Repair scope and query-layout cost}
\label{sec:results}
\subsection{Repair scope changes draft profitability}
\label{sec:economics}
The seven-arm campaign compares repair scope and drafting on the same 12-input panel. All seven arms qualify 96/96 complete endpoints per family under the declared reference. Figure~\ref{fig:cost} separates local representation cost, global serialization cost and the draft increment under all three repair organizations. These estimates retain the original two-session protocol, separately from E2.

\begin{figure*}[t]\centering
\resizebox{\textwidth}{!}{\begin{tikzpicture}
\begin{groupplot}[group style={group size=3 by 1,horizontal sep=.70cm},
width=.305\textwidth,height=5.6cm,ymin=.5,ymax=3.5,ytick={1,2,3},
tick label style={font=\footnotesize},title style={font=\small,yshift=12pt},label style={font=\footnotesize,align=center},
axis lines*=left,axis line style={aux,line width=.45pt},tick style={aux},grid=major,grid style={gray!13}]
\nextgroupplot[title={\textbf{A} Local representation},xlabel={$T_{C_{\PQ}}/T_{C_{\QB}}$\\$>1$: QB faster},xmin=.995,xmax=1.09,xtick={1,1.04,1.08},yticklabels={Llama,SmolLM2,Qwen3}]
\addplot[aux,dashed] coordinates {(1,.5)(1,3.5)};
\addplot[only marks,mark=diamond*,mark size=2.7pt,qualified] coordinates {(1.043676000000,3)};
\node[font=\scriptsize,anchor=south] at (axis cs:1.043676000000,3.12){1.044};
\addplot[only marks,mark=diamond*,mark size=2.7pt,qualified] coordinates {(1.068111000000,2)};
\node[font=\scriptsize,anchor=south] at (axis cs:1.068111000000,2.12){1.068};
\addplot[only marks,mark=diamond*,mark size=2.7pt,qualified] coordinates {(1.021848000000,1)};
\node[font=\scriptsize,anchor=south] at (axis cs:1.021848000000,1.12){1.022};
\nextgroupplot[title={\textbf{B} Global serialization},xlabel={$T_{C_{\GS}}/T_{C_{\QB}}$\\$>1$: QB faster},xmin=.95,xmax=2.55,xtick={1,1.5,2,2.5},yticklabels={,,}]
\addplot[aux,dashed] coordinates {(1,.5)(1,3.5)};
\addplot[only marks,mark=diamond*,mark size=2.7pt,qualified] coordinates {(2.143525000000,3)};
\node[font=\scriptsize,anchor=south] at (axis cs:2.143525000000,3.12){2.144};
\addplot[only marks,mark=diamond*,mark size=2.7pt,qualified] coordinates {(1.462530000000,2)};
\node[font=\scriptsize,anchor=south] at (axis cs:1.462530000000,2.12){1.463};
\addplot[only marks,mark=diamond*,mark size=2.7pt,qualified] coordinates {(2.175941000000,1)};
\node[font=\scriptsize,anchor=south] at (axis cs:2.175941000000,1.12){2.176};
\nextgroupplot[title={\textbf{C} Draft value by repair},xlabel={$D_L=T_{M_L}/T_{C_L}$\\$>1$: drafting faster},xmin=.53,xmax=1.27,xtick={.6,.8,1,1.2},yticklabels={,,},legend style={at={(.5,1.01)},anchor=south,draw=none,font=\scriptsize,fill=none},legend columns=3]
\addplot[aux,dashed,forget plot] coordinates {(1,.5)(1,3.5)};
\addplot[only marks,mark=square*,mark size=2.4pt,aux] coordinates {(0.709722260013,3.23) (0.592109590630,2.23) (0.724338602453,1.23)};
\addlegendentry{GS}
\addplot[only marks,mark=o,mark size=2.4pt,target] coordinates {(1.125118952936,3.00) (0.616360145288,2.00) (1.107396133613,1.00)};
\addlegendentry{PQ}
\addplot[only marks,mark=diamond*,mark size=2.4pt,qualified] coordinates {(1.143749000000,2.77) (0.633948000000,1.77) (1.116620000000,0.77)};
\addlegendentry{QB}
\end{groupplot}
\end{tikzpicture}
}
\caption{Matched repair and drafting choices in the original seven-arm campaign, separate from E2. A/B: speculative PQ and GS costs relative to QB. C: draft value under each repair organization, $D_L=T_{M_L}/T_{C_L}$; values above one favor drafting. Qwen3 and Llama cross parity under both local repairs, while SmolLM2 remains below it. GS/PQ values in C are reconstructed by Equation~\ref{eq:draft} from the common-panel ratios. All points are descriptive estimates; exact source ratios and every regression are retained in \sref{supp:cost}.}
\label{fig:cost}
\end{figure*}

Its speculative GS/QB ratios are 2.144, 1.463 and 2.176; every task favors QB. These differences reflect the value of constraining attention locally while retaining batched projection and feed-forward work. The PQ/QB ratios, 1.044, 1.068 and 1.022, measure the representation choice within local repair. These common-panel estimates distinguish repair scope from layout cost; E2 supplies a separately scheduled PQ/QB comparison.

A retrospective audit of all 2,016 formal responses forms 288 matched seven-arm audit groups. All 1,728 within-path GS/PQ/QB pairings agree in recorded proposals, acceptance, logical rows, rollback and caller-width histograms; complete native token IDs and recorded EOG flags match the paired $R$ endpoints. GS changes physical batching as intended; PQ/QB agree in every workload field. This holds recorded logical work fixed without apportioning operator time (\sref{supp:sevenarmwork}).

Adding the draft to QB is profitable for Qwen3 and Llama but not SmolLM2: $M_{\QB}/C_{\QB}$ is 1.144, 0.634 and 1.117. SmolLM2 takes about 1.58 times its own draft-free request duration despite QB being cheaper than PQ. Across the common seven-arm panel, draft value under another layout $L$ follows
\begin{equation}
D_L=\frac{S_{M_L/M_{\QB}}}{S_{C_L/C_{\QB}}}\,D_{\QB},\qquad D_{\QB}=S_{M_{\QB}/C_{\QB}}.
\label{eq:draft}
\end{equation}
The derived GS/PQ/QB values are 0.710/1.125/1.144 for Qwen3, 0.592/0.616/0.634 for SmolLM2 and 0.724/1.107/1.117 for Llama. Thus the repair policy changes the sign of the observed draft decision for two families. Both PQ and QB cross the profitability boundary; their direct comparison measures QB's additional layout benefit.

The relative change in draft value summarizes this interaction:
\begin{equation}
\Gamma_{\QB\leftarrow\GS}=\frac{D_{\QB}}{D_{\GS}}
=\frac{S_{C_{\GS}/C_{\QB}}}{S_{M_{\GS}/M_{\QB}}}.
\label{eq:repairinteraction}
\end{equation}
It is 1.612, 1.071 and 1.542 for the three families, respectively. These are retrospective ratio-of-ratios summaries from the reported common-panel estimates, with the original weights and timer. They quantify how repair policy changes the draft increment at complete-request level. The draft's incremental cost must therefore be evaluated under the repair policy being considered.

Reference-preserving batching also has value without drafting. Target-only QB lowers cost relative to the specified serial reference: $R/M_{\QB}$ is 1.334051, 1.369890 and 1.426161, with every task favoring QB on this panel. For SmolLM2, the measured choice is therefore to retain target-only QB batching and omit the draft. These complete-request gains include prompt prefill; the target-only continuation uses a common one-query path. They establish neither a single-token decoding speedup nor a measured decomposition by phase.

A historical fixed-UB4 experiment provides a related counterexample: increasing $k$ from 3 to 7 preserved outputs but reduced request speed to 0.890 of the smaller-draft path. Separate work captures recorded fewer target callers and more physical micrographs. Acceptance and caller reduction alone therefore cannot establish request acceleration. The width failures and negative cost results that constrain this conclusion remain in \sref{supp:implementation}.

\subsection{Draft decisions across the six E2 sessions}
\label{sec:e2draft}
E2 records target-only and speculative requests under each local layout. We retrospectively pair its 864 non-warmup local-arm measurements into 216 family/task/session quartets: $M_{\PQ}$, $M_{\QB}$, $C_{\PQ}$ and $C_{\QB}$. The contrast $D_L=S_{M_L/C_L}$ uses the weights in Equation~\ref{eq:ratio}.

\begin{table}[t]\centering\small
\caption{Retrospective E2 draft values on the fixed 12-input panel: $D_L=T_{M_L}/T_{C_L}>1$ favors drafting. QB intervals are exploratory pointwise 95\% $t_5$ intervals over six session log means. E2 has no GS arm.}\label{tab:e2draft}
\begin{tabular}{@{}lrrr@{}}\toprule
Family & $D_{\mathrm{PQ}}$ & $D_{\mathrm{QB}}$ & QB interval\\\midrule
Qwen3 & 1.072 & 1.084 & [1.059, 1.109]\\
SmolLM2 & 0.576 & 0.597 & [0.582, 0.612]\\
Llama & 1.091 & 1.105 & [1.097, 1.113]\\
\bottomrule\end{tabular}
\end{table}


Table~\ref{tab:e2draft} retains all three local-policy directions: drafting benefits Qwen3 and Llama, while SmolLM2 favors target-only execution. Every one of the six session estimates has the corresponding direction under both PQ and QB. Under QB, drafting reduces complete-request latency by 7.72\% and 9.51\% for Qwen3 and Llama; it increases SmolLM2 latency by 67.52\%. The magnitudes differ from the original two-session campaign, while the policy decisions agree.

The exploratory pointwise $t_5$ intervals use sessions as sampling units; the registered PQ/QB analysis is unchanged. Pairing separately collected target-only and speculative blocks does not remove within-session drift. Without a GS arm, E2 checks the local draft decision rather than independently replicating the full repair-scope interaction. Complete estimates and the PQ-to-QB change in draft value appear in \sref{supp:e2draft}.

\subsection{QB's independent implementation benefit}
\label{sec:e2}
The registered E2 comparison isolates representation cost within local repair. It completes 1,620 requests in 324 fresh workers, including 324 warmups and 432 serial-reference measurements. There are no request failures, output mismatches, workload errors or heavy-observer flags. The six-session speculative ratios $C_{\PQ}/C_{\QB}$ are 1.031614, 1.065803 and 1.029114, corresponding to 3.06\%, 6.17\% and 2.83\% lower QB request latency. Figure~\ref{fig:e2} shows both paths and their session-level uncertainty.

Table~\ref{tab:e2absolute} gives absolute costs from the source-verified requests, using the same task/session log weights.
\begin{table}[t]\centering\small
\caption{Absolute E2 request costs (ms), with equal task and session weights in the geometric mean. Registered warmups are excluded; the paid interval includes submission, full-response parsing and endpoint validation.}\label{tab:e2absolute}
\begin{tabular}{lrrrr}
\toprule
Family & $M_{\mathrm{PQ}}$ & $M_{\mathrm{QB}}$ & $C_{\mathrm{PQ}}$ & $C_{\mathrm{QB}}$ \\
\midrule
Qwen3 & 2130.4 & 2087.5 & 1987.2 & 1926.3 \\
SmolLM2 & 713.1 & 693.8 & 1238.8 & 1162.3 \\
Llama & 2574.7 & 2535.1 & 2360.9 & 2294.2 \\
\bottomrule
\end{tabular}

\end{table}

\begin{figure*}[t]\centering
\begin{tikzpicture}[font=\footnotesize]
\begin{axis}[name=sess,scale only axis,width=13.4cm,height=4.35cm,
xmin=.998,xmax=1.078,xtick={1,1.02,1.04,1.06},xticklabels={1.00,1.02,1.04,1.06},
ymin=.55,ymax=6.55,ytick={6,5,4,3,2,1},yticklabels={Qwen3 M,Qwen3 C,SmolLM2 M,SmolLM2 C,Llama M,Llama C},
xlabel={PQ / QB request-time ratio ($>1$: QB faster)},xlabel style={font=\footnotesize},tick label style={font=\footnotesize},
axis lines*=left,axis line style={aux},tick style={aux},grid=major,grid style={gray!13}]
\draw[aux,dashed] (axis cs:1,.55)--(axis cs:1,6.55);
\draw[draft,densely dashed] (axis cs:1.02,.55)--(axis cs:1.02,6.55);
\addplot[only marks,mark=*,mark size=1.45pt,mark options={draw=aux,fill=white,line width=.5pt},color=aux] coordinates {(1.022241155,6.2) (1.003825879,6.2) (1.024056173,6.2) (1.031462696,6.2) (1.020617513,6.2) (1.021379783,6.2)};
\addplot[only marks,mark=*,mark size=1.45pt,mark options={draw=aux,fill=white,line width=.5pt},color=aux] coordinates {(1.018417206,5.2) (1.034686907,5.2) (1.029528424,5.2) (1.027889260,5.2) (1.044209114,5.2) (1.035132037,5.2)};
\addplot[only marks,mark=*,mark size=1.45pt,mark options={draw=aux,fill=white,line width=.5pt},color=aux] coordinates {(1.021523343,4.2) (1.042811401,4.2) (1.026887696,4.2) (1.027267920,4.2) (1.014728654,4.2) (1.033614105,4.2)};
\addplot[only marks,mark=*,mark size=1.45pt,mark options={draw=aux,fill=white,line width=.5pt},color=aux] coordinates {(1.062659839,3.2) (1.069420131,3.2) (1.066117550,3.2) (1.070418529,3.2) (1.063618892,3.2) (1.062610884,3.2)};
\addplot[only marks,mark=*,mark size=1.45pt,mark options={draw=aux,fill=white,line width=.5pt},color=aux] coordinates {(1.011619881,2.2) (1.003629727,2.2) (1.026339325,2.2) (1.016926078,2.2) (1.023407939,2.2) (1.011808397,2.2)};
\addplot[only marks,mark=*,mark size=1.45pt,mark options={draw=aux,fill=white,line width=.5pt},color=aux] coordinates {(1.013167868,1.2) (1.050133956,1.2) (1.021893776,1.2) (1.036275997,1.2) (1.032350688,1.2) (1.021280772,1.2)};
% Asymmetric percentile intervals and point estimates are generated from the frozen table.
\draw[target,line width=1pt] (axis cs:1.013022800003,6)--(axis cs:1.026553489035,6);
\draw[target,line width=.65pt] (axis cs:1.013022800003,5.925)--(axis cs:1.013022800003,6.075);
\draw[target,line width=.65pt] (axis cs:1.026553489035,5.925)--(axis cs:1.026553489035,6.075);
\addplot[only marks,mark=o,mark size=2.4pt,target] coordinates {(1.020563161899,6)};
\draw[qualified,line width=1pt] (axis cs:1.025392736138,5)--(axis cs:1.037925690306,5);
\draw[qualified,line width=.65pt] (axis cs:1.025392736138,4.925)--(axis cs:1.025392736138,5.075);
\draw[qualified,line width=.65pt] (axis cs:1.037925690306,4.925)--(axis cs:1.037925690306,5.075);
\addplot[only marks,mark=square*,mark size=2.4pt,qualified] coordinates {(1.031613690048,5)};
\draw[target,line width=1pt] (axis cs:1.021009792824,4)--(axis cs:1.035016679387,4);
\draw[target,line width=.65pt] (axis cs:1.021009792824,3.925)--(axis cs:1.021009792824,4.075);
\draw[target,line width=.65pt] (axis cs:1.035016679387,3.925)--(axis cs:1.035016679387,4.075);
\addplot[only marks,mark=o,mark size=2.4pt,target] coordinates {(1.027767425686,4)};
\draw[qualified,line width=1pt] (axis cs:1.063387048182,3)--(axis cs:1.068398798323,3);
\draw[qualified,line width=.65pt] (axis cs:1.063387048182,2.925)--(axis cs:1.063387048182,3.075);
\draw[qualified,line width=.65pt] (axis cs:1.068398798323,2.925)--(axis cs:1.068398798323,3.075);
\addplot[only marks,mark=square*,mark size=2.4pt,qualified] coordinates {(1.065803002589,3)};
\draw[target,line width=1pt] (axis cs:1.009627059351,2)--(axis cs:1.021621860041,2);
\draw[target,line width=.65pt] (axis cs:1.009627059351,1.925)--(axis cs:1.009627059351,2.075);
\draw[target,line width=.65pt] (axis cs:1.021621860041,1.925)--(axis cs:1.021621860041,2.075);
\addplot[only marks,mark=o,mark size=2.4pt,target] coordinates {(1.015593036909,2)};
\draw[qualified,line width=1pt] (axis cs:1.020125595877,1)--(axis cs:1.039148610748,1);
\draw[qualified,line width=.65pt] (axis cs:1.020125595877,0.925)--(axis cs:1.020125595877,1.075);
\draw[qualified,line width=.65pt] (axis cs:1.039148610748,0.925)--(axis cs:1.039148610748,1.075);
\addplot[only marks,mark=square*,mark size=2.4pt,qualified] coordinates {(1.029113522487,1)};
\end{axis}
\end{tikzpicture}

\caption{Independent-block E2 replication. Small marks are the six session estimates; large markers and horizontal bars are the equal-weight estimate and its registered 95\% six-session percentile interval. Dashed lines mark parity and the 1.02 practical threshold. $M$ is target-only and $C$ speculative. The fixed exposed panel and single host limit statistical scope.}
\label{fig:e2}
\end{figure*}

All 36 family/path/session contrasts favor QB; 9 of the 72 family/path/task aggregates regress. The primary pointwise speculative intervals are [1.025393,1.037926], [1.063387,1.068399] and [1.020126,1.039149], each above the predeclared 1.02 threshold under the registered percentile method. Target-only effects are smaller: Llama's estimate is 1.015593, below the threshold, while Qwen3's interval overlaps it.

The secondary $t_5$ intervals remain above 1.02 for speculative Qwen3 and SmolLM2. Llama's interval is [1.015387,1.043026], above parity but crossing 1.02; its six leave-one-session-out estimates range from 1.024960 to 1.032333. QB's advantage is stable in these checks. Whether the Llama effect exceeds the practical threshold depends on the interval method. \sref{supp:cuda} reports both analyses and all task regressions.

E2 establishes endpoint and recorded-counter correspondence across all 216 speculative PQ/QB task/session pairs: acceptance, rollback and logical-workload counters agree. This supports an implementation-cost comparison under matched recorded work. The original two-session estimates and their target-only direction reversals are reported in \sref{supp:cost}; E2's recorded-counter scope is detailed in \sref{supp:cuda}.

The recovered E2 records do not report actual CUDA Graph capture, replay or fallback counts. The request-cost estimate therefore identifies the implementation effect but does not isolate those mechanisms. The controlled QB/QH replay study in Section~\ref{sec:mappingcosts} has its own execution evidence and comparator.


\subsection{Same-input correspondence across layers and phases}
\label{sec:numeric}
The initial 24-context replay established PQ/QB correspondence at sampled layer-0 attention boundaries. E1 expands the test to three layer positions, prefill and continuation, and $q=1,2,3,4$ in each family. All 72 contexts pass observation, native-MQ own-path and repeatability checks. PQ and QB are bitwise equal at QK, softmax and AV in every context, with zero maximum error and ULP distance (Figure~\ref{fig:numeric}).

\begin{figure}[t]\centering
\resizebox{\columnwidth}{!}{\begin{tikzpicture}[x=1cm,y=1cm,font=\footnotesize]
\node[anchor=west,font=\bfseries\small] at (0,2.65) {Bitwise-equal contexts / 72};
\foreach \x/\t in {3.30/QK,4.95/Softmax,6.60/AV}\node at (\x,2.15) {\t};
\foreach \y/\t in {1.55/MQ vs PQ,.85/MQ vs QB,.15/PQ vs QB}\node[anchor=west] at (0,\y) {\t};
\foreach \y in {1.55,.85}{\foreach \x in {3.30,4.95,6.60}{
 \draw[draft,line width=.65pt] (\x-.64,\y-.25) rectangle (\x+.64,\y+.25);
 \node[text=draft] at (\x,\y) {18 / 72};}}
\foreach \x in {3.30,4.95,6.60}{\draw[qualified,line width=.65pt,fill=qualified!8] (\x-.64,-.10) rectangle (\x+.64,.40);\node[text=qualified] at (\x,.15) {72 / 72};}
\node[anchor=west,text=aux,font=\scriptsize] at (0,-.55) {18 single-query controls; 54 multi-query contexts};
\end{tikzpicture}
}
\caption{Completed 72-context CUDA replay. The 18 agreements between MQ and either repair are exactly the $q=1$ common-path controls. The 54 multi-query cases first differ at QK. All PQ/QB stage outputs agree bitwise.}
\label{fig:numeric}
\end{figure}

The sampled target layers are $0/18/35$ for Qwen3, $0/12/23$ for SmolLM2 and $0/16/31$ for Llama. Qwen3 and Llama use head dimension 128 with 32 query heads and 8 KV heads; SmolLM2 uses dimension 64 with 32/32 heads. Thus the observed correspondence spans both ordinary and grouped-query attention and more than one head dimension.

MQ agrees with either repaired path only in the 18 single-query controls. Across the other 54 contexts, its first difference is at QK; maximum finite absolute differences are 0.078583 at QK, 0.000917 at softmax and 0.010433 at AV. These are operator differences on identical operands, not model-logit bounds.

All snapshots have tensor key extent 256, including masked padding. The maximum number of visible keys is 109, 130 and 127 by family; prefill captures include early query positions with only one to four visible keys. The experiment covers layers, phases and query widths at these short visible contexts, using one exposed prompt per family. Complete coordinates and operator comparisons appear in \sref{supp:cuda}.

